\documentclass[article]{jss}

%% -- Extra packages not in jss.cls -----------------------------------------------
\usepackage{amsmath,amssymb}
\usepackage{booktabs,array,longtable,multirow}
\usepackage{graphicx}
\usepackage{enumitem}
\usepackage{subcaption}
\usepackage{tcolorbox}

%% -- Code box styles (override jss.cls defaults) ---------------------------------
\definecolor{codebackground}{RGB}{248,248,248}
\definecolor{outputbackground}{RGB}{240,248,240}
\tcbset{
  codebox/.style={
    colback=codebackground, colframe=gray!50,
    left=4pt, right=4pt, top=2pt, bottom=2pt,
    fontupper=\small\ttfamily, boxrule=0.5pt, arc=2pt
  },
  outputbox/.style={
    colback=outputbackground, colframe=gray!40,
    left=4pt, right=4pt, top=2pt, bottom=2pt,
    fontupper=\small\ttfamily, boxrule=0.5pt, arc=2pt
  }
}

%% -- Title -----------------------------------------------------------------------
\title{\pkg{lepidoptera}: A Domain-Agnostic Engine for Causal Simulation
  with Mixed Variable Types, Vine Copula Dependence,
  and Counterfactual Propagation}

\titleshort{\pkg{lepidoptera}: Causal Simulation Engine}

\author{Danilo A.~Sarti\\
        Universidade Estadual de Campinas}

\authorshort{Sarti}

\Keywords{structural causal model, directed acyclic graph, mixed variable types,
  vine copula, counterfactual simulation, partial propagation,
  \proglang{R}, \proglang{Python}}

\Plainkeywords{structural causal model, directed acyclic graph, mixed variable
  types, vine copula, counterfactual simulation, partial propagation, R, Python}

%% ══════════════════════════════════════════════════════════════════════════════
\begin{document}

\maketitle

\begin{abstract}
\noindent
We introduce \pkg{lepidoptera}, an open-source \proglang{R} and \proglang{Python}
package providing a modular, domain-agnostic engine for simulation from structural
causal models (SCMs) with heterogeneous variable types. The package makes four
contributions not available in combination in any existing software. First, a
\emph{typed mechanism} architecture separates the causal graph from the generating
rules at each node, exposing each rule as an inspectable, serialisable
\class{Mechanism} object that can be replaced in isolation for any intervention or
sensitivity analysis. Second, a \emph{pre-built mechanism library} covers all nine
parent-type / child-type combinations (continuous, categorical, ordinal $\times$
continuous, categorical, ordinal), including linear Gaussian, multinomial logit,
ordered logit, group-effect, and categorical transition models, removing the need
for users to implement type dispatch from scratch. Third, a \emph{vine copula
mechanism} learns the full joint dependence structure from observed data via
pair-copula constructions and the Rosenblatt conditional-simulation algorithm,
supporting over twenty copula families from \pkg{rvinecopulib}
\citep{nagler2022rvinecopulib}. Fourth, \emph{partial propagation} operationalises
individual-level counterfactual queries by freezing the random draws of unaffected
nodes and replaying only the downstream consequences of an intervention, cleanly
separating the ITE estimand from the population-level ATE without requiring
additional model assumptions. An AI-native extension layer allows any pre-fitted
predictive model --- random forests, gradient-boosted trees, neural networks --- to
serve as a generative mechanism. Simulation studies validate type-combination
coverage, copula family fidelity, and counterfactual consistency. Real-data
illustrations span clinical trial simulation, hospital emergency-ward digital
twinning, and epidemiological mediation analysis. \pkg{lepidoptera} is available
under the MIT licence at \url{https://github.com/daniloasarti/lepidoptera}.

\medskip\noindent
\textbf{Keywords:} structural causal model, directed acyclic graph, mixed variable
types, vine copula, counterfactual simulation, partial propagation,
\proglang{R}, \proglang{Python}.
\end{abstract}

\tableofcontents
\newpage

%% ══════════════════════════════════════════════════════════════════════════════
\section{Introduction}
\label{sec:intro}

\subsection{Motivation}

Consider three scenarios that regularly arise in applied statistics.

\begin{example}[Clinical trial planning]
A biostatistician designs a Phase~II oncology trial. The patient population includes
continuous covariates (age, tumour size, biomarker level), an ordinal baseline
performance score (ECOG 0--4), and a categorical treatment group. The primary
endpoint is a binary indicator of response, confounded by performance score. The
analyst wants to (i) generate a realistic synthetic population that honours the
observed joint distribution of these heterogeneous covariates, (ii) apply an
intervention that fixes the treatment assignment for all patients (\emph{population
policy}), and (iii) ask, for each patient who actually failed, what their outcome
\emph{would have been} under a different dosing regimen (\emph{individual
counterfactual}).
\end{example}

\begin{example}[Emergency-ward digital twin]
A hospital data scientist builds a simulation engine to stress-test staffing
protocols. The model includes continuous inter-arrival times, ordinal triage
categories (Priority 1--3), continuous length-of-stay, and a categorical discharge
destination. She wishes to intervene on triage protocols and measure the
downstream effect on boarding times for individual patients, not just the
population average.
\end{example}

\begin{example}[Epidemiological mediation]
An epidemiologist wants to decompose the effect of income on cardiovascular
disease risk into a direct path and a path mediated through dietary quality and
smoking status. Income is continuous; smoking is categorical; diet is continuous;
CVD risk is binary. Estimating natural direct and indirect effects requires
simulating counterfactual outcomes under joint interventions on both mediators,
for each individual unit.
\end{example}

All three scenarios share a common computational need: simulate from a DAG with
mixed variable types, intervene on one or more nodes, and distinguish between a
population-level effect and an individual-level (counterfactual) effect. Despite
the prevalence of this need, no single existing \proglang{R} or \proglang{Python}
package satisfies all three requirements. Existing tools either restrict mechanisms
to a single variable type, do not expose the mechanism as a replaceable object, or
conflate population-level and individual-level propagation.

We present \pkg{lepidoptera}, an open-source package that fills this gap. The name
recalls the modular, compositional logic of biological taxonomy: just as a
lepidopterist builds an identification key from composable, observable features,
the package builds a data-generating process from composable, observable
mechanisms---each typed, replaceable, and independently auditable.

\subsection{Overview of contributions}

The central design decision in \pkg{lepidoptera} is the strict separation of
\emph{graph} (which variables exist and who depends on whom) from
\emph{mechanism} (how each variable is generated given its parents). This
separation, grounded in Pearl's \emph{modularity assumption} for SCMs
\citep{pearl2009causality}, has consequences that propagate throughout the package:

\begin{itemize}[leftmargin=1.8em]
  \item \textbf{Type-complete mechanism library.} All nine
        (parent type, child type) combinations are covered by pre-built factories.
        Users specify parameters; the package handles type coercion, domain
        validation, and the correct sampling distribution.
  \item \textbf{Vine copula mechanism.} A single call to
        \fct{copula\_mechanism} fits a pair-copula construction to observed data
        and registers a conditional-simulation rule at the specified node, using
        the Rosenblatt transform \citep{rosenblatt1952remarks} for exact
        inversion.
  \item \textbf{Partial propagation.} \fct{propagate(..., mode = "partial")}
        freezes all unaffected random draws and recomputes only the downstream
        subgraph, enabling ITE estimation from the same synthetic population as
        the ATE.
  \item \textbf{AI-native integration.} Any \proglang{R} or \proglang{Python}
        callable can serve as a mechanism via \fct{model\_mechanism},
        \fct{learned\_mechanism}, or \fct{llm\_mechanism}, without any adapter
        code.
  \item \textbf{Cross-language portability.} Systems serialise to JSON; a model
        defined in \proglang{R} can be imported and re-simulated in
        \proglang{Python} and vice versa.
\end{itemize}

\subsection{Paper organisation}

Section~\ref{sec:related} surveys the related literature and positions
\pkg{lepidoptera} against the closest existing tools. Section~\ref{sec:framework}
gives formal definitions. Section~\ref{sec:mechanisms} describes the mechanism
library with code examples. Section~\ref{sec:simulation} covers the simulation
and propagation algorithms. Section~\ref{sec:counterfactual} formalises
counterfactual propagation and proves the consistency result. Section~\ref{sec:copula}
presents the vine copula mechanism. Section~\ref{sec:ai} describes the AI-native
layer. Section~\ref{sec:sims} reports simulation studies. Section~\ref{sec:applications}
presents three worked applications. Section~\ref{sec:discussion} discusses
limitations and future directions.

%% ══════════════════════════════════════════════════════════════════════════════
\section{Related Work}
\label{sec:related}

\subsection{Causal simulation in \proglang{R}}

\paragraph{\pkg{simcausal}.}
\citet{sofrygin2017simcausal} introduced \pkg{simcausal}, which provides a
formula-based domain-specific language for specifying non-parametric SEMs for
longitudinal data with time-varying treatments, censoring, and monitoring
processes. The package computes causal quantities such as treatment-specific means
and coefficients from marginal structural models \citep{robins2000marginal}, and
has been widely used to benchmark targeted learning estimators
\citep{luedtke2016tloss, laan2018targeted}. Its design, however, is oriented
around a temporal node index and a shared \proglang{R} environment evaluator,
making it cumbersome for cross-sectional DAGs with heterogeneous mechanisms. There
is no typed mechanism object, no vine copula, and no native support for
individual-level (counterfactual) queries.

\paragraph{\pkg{simstudy}.}
\citet{goldfeld2020simstudy} provide \pkg{simstudy}, an expressive \proglang{R}
package for study-design simulation supporting binary, Poisson, negative-binomial,
gamma, and survival outcomes. Its formula interface is convenient for power
calculations, but the package lacks a causal graph layer; intervention propagation
must be implemented manually by the user.

\paragraph{\pkg{simDAG}.}
\citet{denz2026simdag} present \pkg{simDAG}, published in the Journal of Statistical
Software (vol.~116, no.~2), which generates data from a DAG using node type
descriptors (\code{node\_gaussian}, \code{node\_binomial}, \code{node\_multinomial},
\code{node\_cox}, etc.) for both cross-sectional and discrete-time longitudinal
settings. \pkg{simDAG} is the closest predecessor to \pkg{lepidoptera} in the
\proglang{R} ecosystem, supporting a broader type palette than \pkg{simcausal} and
covering mixed-type DAGs. It differs from \pkg{lepidoptera} in three key respects:
(i) mechanisms are node-type strings rather than first-class objects; (ii) there is
no vine copula mechanism for data-driven dependence modelling; and (iii)
counterfactual (individual-level) propagation is not natively supported.

\paragraph{\pkg{bnlearn}.}
\citet{scutari2010bnlearn} provide \pkg{bnlearn} for Bayesian network structure
learning and inference. Simulation is restricted to linear Gaussian or discrete
conditional probability table mechanisms; ordinal variables, vine copulas, and
partial propagation are absent.

\paragraph{\pkg{R6causal}.}
\citet{tikka2024r6causal} provide \pkg{R6causal}, an R6-class implementation of
SCMs with functions for data generation, do-calculus identification (via
\pkg{causaleffect} and \pkg{dosearch}), and counterfactual simulation. The package
introduces a \class{ParallelWorld} class for discrete counterfactual inference.
Its scope is narrower than \pkg{lepidoptera}: mechanisms are restricted to
parametric forms, support for continuous counterfactuals requires additive error
structures, and mixed ordinal/categorical types lack dedicated mechanism factories.

\paragraph{\pkg{dagitty}.}
\citet{textor2016dagitty} provide a web and \proglang{R} interface for DAG
analysis (d-separation, adjustment sets, equivalent models). The package does not
generate data; \pkg{lepidoptera} is complementary, not competing.

\paragraph{\pkg{lavaan} and \pkg{OpenMx}.}
\citet{rosseel2012lavaan} and \citet{neale2016openmx} implement structural equation
models in \proglang{R}. Both are restricted to linear Gaussian mechanisms and are
better positioned as estimation tools than simulation engines.

\subsection{Causal simulation in \proglang{Python}}

\paragraph{\pkg{DagSim}.}
\citet{alhajj2023dagsim} introduce \pkg{DagSim}, a \proglang{Python} framework
that removes restrictions on variable types and functional forms by routing all
generating logic through user-supplied functions specified in YAML. The framework
is motivated by the observation that existing DAG-based simulation packages are
still confined to settings with relatively simple variable types and functional
forms \citep{alhajj2023dagsim}. \pkg{lepidoptera}'s Python implementation follows
a similar philosophy but adds typed mechanism validation, a pre-built mechanism
library, vine copula integration, and partial propagation.

\paragraph{\pkg{causalnex} and \pkg{DoWhy}.}
\citet{beaumont2021causalnex} and \citet{sharma2020dowhy} focus on structure
learning and causal effect estimation, not on generative simulation.

\subsection{Counterfactual simulation}

\citet{karvanen2024simcf} present an algorithm for exact counterfactual simulation
in parametric SCMs where the structural equations are monotonic or additive in
their error terms. Their method handles background variables that affect multiple
observed variables, a setting where prior algorithms failed. The method is more
general than \pkg{lepidoptera}'s partial propagation in that it handles
non-Gaussian error structures explicitly, but it requires fully parametric models
and does not target the mixed-type simulation workflow.

The \pkg{cfid} package \citep{tikka2022cfid} addresses \emph{identification} of
counterfactual queries via the $\mathrm{ID}^*$ and $\mathrm{IDC}^*$ algorithms
\citep{shpitser2007id}; it does not simulate data.

\subsection{Synthetic data and copula models}

\citet{nowok2016synthpop} introduce \pkg{synthpop} for synthetic data release via
sequential regression imputation, optionally guided by a user-supplied DAG ordering.
The focus is on disclosure-limited synthetic data rather than causal simulation;
mechanisms are regression models fitted to a real dataset rather than user-specified.

For copula modelling, \pkg{rvinecopulib} \citep{nagler2022rvinecopulib} and
\pkg{VineCopula} \citep{nagler2022rvinecopulib} are the state of the art in
\proglang{R}. \pkg{lepidoptera} wraps \pkg{rvinecopulib} rather than
reimplementing it, inheriting its full family set and C++ performance.
\pkg{pyvinecopulib} \citep{pyvinecopulib2024} provides the same backend for the
\proglang{Python} implementation.

\subsection{Summary gap analysis}

Table~\ref{tab:comparison} places \pkg{lepidoptera} in the feature matrix of the
closest existing tools. The combination of typed mechanisms, full mixed-type
coverage, vine copula support, and individual-level counterfactual propagation is
unique to \pkg{lepidoptera}.

\begin{table}[ht]
\centering
\caption{Feature comparison of causal simulation tools (October 2026). ``Mixed
types'' = native support for continuous, categorical, and ordinal in any
parent-child combination. ``Copula'' = vine copula mechanism. ``Counterfactual''
= individual-level partial propagation. ``Typed mech.'' = mechanisms as
first-class objects. ``R/Py'' = availability in both languages.}
\label{tab:comparison}
\small\setlength{\tabcolsep}{5pt}
\begin{tabular}{lcccccr}
\toprule
Package & Language & Mixed types & Copula & Counterfactual & Typed mech. \\
\midrule
\pkg{simcausal} \citep{sofrygin2017simcausal} & R & Partial & No & No & No \\
\pkg{simstudy}  \citep{goldfeld2020simstudy}  & R & Partial & No & No & No \\
\pkg{simDAG}    \citep{denz2026simdag}         & R & Yes & No & No & No \\
\pkg{bnlearn}   \citep{scutari2010bnlearn}     & R & Partial & No & No & No \\
\pkg{R6causal}  \citep{tikka2024r6causal}      & R & Partial & No & Partial & No \\
\pkg{lavaan}    \citep{rosseel2012lavaan}       & R & No & No & No & No \\
\pkg{dagitty}   \citep{textor2016dagitty}       & R & No (viz) & No & No & No \\
\pkg{DagSim}    \citep{alhajj2023dagsim}        & Python & Yes & No & No & No \\
\pkg{causalnex} \citep{beaumont2021causalnex}  & Python & Partial & No & No & No \\
\pkg{lepidoptera} (this paper)                 & R+Python & \textbf{Yes} & \textbf{Yes} & \textbf{Yes} & \textbf{Yes} \\
\bottomrule
\end{tabular}
\end{table}

%% ══════════════════════════════════════════════════════════════════════════════
\section{The \pkg{lepidoptera} Framework}
\label{sec:framework}

\subsection{Variable types and domains}
\label{sec:types}

\begin{definition}[Variable type and domain]
\label{def:type}
A \emph{feature} $V$ has a \emph{type} $\tau(V)$ drawn from the closed set
$\mathcal{T} = \{\mathtt{continuous},\, \mathtt{categorical},\, \mathtt{ordinal}\}$
and a \emph{domain} $\mathcal{D}(V)$:
\begin{itemize}[nosep]
  \item $\tau(V) = \mathtt{continuous}$: $\mathcal{D}(V) = [a,b] \subseteq
        \overline{\mathbb{R}}$, possibly unbounded.
  \item $\tau(V) = \mathtt{categorical}$: $\mathcal{D}(V) = \mathcal{L}$, a
        finite set with no ordering.
  \item $\tau(V) = \mathtt{ordinal}$: $\mathcal{D}(V) = (\mathcal{L},\prec)$,
        a finite totally ordered set.
\end{itemize}
\end{definition}

\subsection{Mechanisms}
\label{sec:mech_def}

\begin{definition}[Mechanism]
\label{def:mech}
A \emph{mechanism} for feature $V_i$ with parent set $\mathbf{Pa}_i \subseteq
\mathcal{V}\setminus\{V_i\}$ is a triple
$\mathcal{M}_i = \bigl(\mathrm{fn}_i,\, \boldsymbol\tau_{\mathbf{Pa}},\, \tau_i\bigr)$
where:
\begin{enumerate}[nosep,leftmargin=2em]
  \item $\mathrm{fn}_i : \prod_{j \in \mathbf{Pa}_i}\mathcal{D}(V_j) \times \Omega
        \to \mathcal{D}(V_i)$ is a measurable function; $\Omega$ is a probability
        space carrying all exogenous noise.
  \item $\boldsymbol\tau_{\mathbf{Pa}} = \bigl(\tau(V_j)\bigr)_{j\in\mathbf{Pa}_i}$
        is the \emph{declared parent type vector}, checked at
        \fct{finalize\_system} time.
  \item $\tau_i = \tau(V_i)$ is the declared output type.
\end{enumerate}
When $\mathbf{Pa}_i = \emptyset$, the mechanism is \emph{marginal}; otherwise it
is \emph{conditional}.
\end{definition}

In \proglang{R}, a \class{mechanism} object is a named list with fields
\code{name}, \code{fn}, \code{parent\_types}, \code{output\_type}, and
\code{spec} (a serialisable parameter record, or \code{NULL} for closures that
carry model objects).

\subsection{Systems}
\label{sec:system_def}

\begin{definition}[System]
\label{def:system}
A \emph{system} $\mathcal{S} = (\mathcal{V}, \mathcal{E}, \{\mathcal{M}_i\}_{i=1}^p)$
consists of:
\begin{enumerate}[nosep,leftmargin=2em]
  \item A finite set of typed features $\mathcal{V}$.
  \item A set of directed edges $\mathcal{E} \subseteq \mathcal{V}\times\mathcal{V}$
        such that the induced graph $G = (\mathcal{V}, \mathcal{E})$ is acyclic.
  \item A mechanism $\mathcal{M}_i$ for each $V_i$, with
        $\boldsymbol\tau_{\mathbf{Pa}}$ consistent with the actual types of
        $\{V_j : (V_j,V_i)\in\mathcal{E}\}$.
\end{enumerate}
\fct{finalize\_system} verifies five invariants: (I1) $G$ is acyclic; (I2) every
$V_i$ has exactly one $\mathcal{M}_i$; (I3) all referenced names exist in
$\mathcal{V}$; (I4) declared parent types match actual types; (I5) no orphaned
mechanisms exist.
\end{definition}

\subsection{Interventions}
\label{sec:intervention_def}

\begin{definition}[Intervention and affected set]
\label{def:intervention}
An \emph{intervention} $\mathcal{I} = (\mathcal{T}, k, \theta)$ specifies:
a \emph{target set} $\mathcal{T}\subseteq\mathcal{V}$, a \emph{kind}
$k \in \{\mathtt{do\_fix},\, \mathtt{do\_distribution},\, \mathtt{shift},\,
\mathtt{mechanism\_override}\}$, and \emph{parameters} $\theta$.

\begin{itemize}[nosep,leftmargin=2em]
  \item $\mathtt{do\_fix}(V_i = v_0)$: $\mathrm{fn}_i \equiv v_0$; all edges
        into $V_i$ are cut.
  \item $\mathtt{do\_distribution}(V_i, \mathcal{M}_i')$: replace $\mathcal{M}_i$
        with a new mechanism $\mathcal{M}_i'$.
  \item $\mathtt{shift}(V_i, \delta)$: $\mathrm{fn}_i' \equiv
        \delta\bigl(\mathrm{fn}_i(\mathbf{pa}_i, \omega)\bigr)$ for a
        measurable $\delta:\mathcal{D}(V_i)\to\mathcal{D}(V_i)$.
  \item $\mathtt{mechanism\_override}(V_i, \mathcal{M}_i'')$: replace
        $\mathcal{M}_i$ with an arbitrary mechanism.
\end{itemize}

Interventions are immutable; they are applied to a \emph{copy}
$\mathcal{S}' = (\mathcal{V}, \mathcal{E}', \{\mathcal{M}_i'\})$ of $\mathcal{S}$,
leaving the original unchanged.

The \emph{affected set} is the transitive downstream closure of $\mathcal{T}$:
\[
  \mathrm{Aff}(\mathcal{T}, G) = \bigl\{V\in\mathcal{V} : \exists\, V_i\in\mathcal{T},\,
  \exists\text{ directed path } V_i \to \cdots \to V \text{ in } G\bigr\}.
\]
\end{definition}

%% ══════════════════════════════════════════════════════════════════════════════
\section{Mechanism Library}
\label{sec:mechanisms}

\subsection{Pre-built factories}
\label{sec:prebuilt}

Table~\ref{tab:mechanisms} lists the nine default mechanism factories, one per
parent-type/child-type combination. All are accessible via
\fct{mechanism\_library\_get}.

\begin{table}[ht]
\centering
\caption{Pre-built mechanism factories in \pkg{lepidoptera}. C = continuous,
K = categorical, O = ordinal. All factories return a \class{mechanism} object.}
\label{tab:mechanisms}
\small
\begin{tabular}{lll}
\toprule
Parent & Child & Factory and model \\
\midrule
C & C & \code{linear\_gaussian}: $Y = \alpha + \boldsymbol\beta^\top\mathbf{X}+\epsilon$,
        $\epsilon\sim\mathcal{N}(0,\sigma^2)$ \\
C & K & \code{multinomial\_logit}: $P(Y\!=\!k\mid\mathbf{X})\propto\exp(\alpha_k+\boldsymbol\beta_k^\top\mathbf{X})$ \\
C & O & \code{ordered\_logit}: proportional-odds \citep{mccullagh1980regression} \\
K & C & \code{group\_effect}: $Y\mid X\!=\!k\sim\mathcal{N}(\mu_k,\sigma_k^2)$ \\
K & K & \code{categorical\_transition}: full $|\mathcal{L}_X|\times|\mathcal{L}_Y|$ conditional table \\
K & O & \code{categorical\_transition} with ordinal output \\
O & C & \code{group\_effect} or \code{linear\_gaussian} on numeric code \\
O & K & \code{categorical\_transition} \\
O & O & \code{categorical\_transition} with optional monotone constraint \\
\midrule
\multicolumn{2}{l}{Any combination} & \code{copula\_mechanism} (Section~\ref{sec:copula}); \\
                                    &  & \code{learned\_mechanism}, \code{model\_mechanism} (Section~\ref{sec:ai}) \\
\bottomrule
\end{tabular}
\end{table}

\paragraph{Linear Gaussian.}
For a continuous child $Y$ with continuous parent vector $\mathbf{X}$:
\begin{equation}
  Y = \alpha + \sum_{j=1}^q \beta_j X_j + \epsilon, \qquad \epsilon\sim\mathcal{N}(0,\sigma^2).
  \label{eq:linear_gaussian}
\end{equation}

\paragraph{Multinomial logit.}
For a categorical child with $K$ unordered levels:
\begin{equation}
  P(Y = k \mid \mathbf{X}) = \frac{\exp(\alpha_k + \boldsymbol\beta_k^\top\mathbf{X})}
  {\sum_{\ell=1}^K \exp(\alpha_\ell + \boldsymbol\beta_\ell^\top\mathbf{X})},
  \label{eq:multinomial}
\end{equation}
with one level designated as the reference ($\alpha_{\mathrm{ref}}=0$,
$\boldsymbol\beta_{\mathrm{ref}}=\mathbf{0}$).

\paragraph{Ordered logit.}
For an ordinal child with levels $\ell_1\prec\cdots\prec\ell_K$
\citep{mccullagh1980regression}:
\begin{equation}
  \log\frac{P(Y\preceq\ell_k\mid\mathbf{X})}{P(Y\succ\ell_k\mid\mathbf{X})}
  = \gamma_k - \boldsymbol\beta^\top\mathbf{X}, \qquad
  \gamma_1 < \cdots < \gamma_{K-1}.
  \label{eq:ordered_logit}
\end{equation}

\paragraph{Group effect.}
For a continuous child with a single categorical or ordinal parent $X$:
\begin{equation}
  Y \mid X = k \;\sim\; \mathcal{N}(\mu_k, \sigma_k^2).
  \label{eq:group_effect}
\end{equation}

\paragraph{Categorical transition.}
For a categorical/ordinal child with a single categorical/ordinal parent, the
mechanism is specified by a full conditional probability table
$\{P(Y = \ell \mid X = k)\}_{k,\ell}$.

\subsection{Code examples}
\label{sec:mech_code}

\paragraph{Example: building a mixed-type system.}

\begin{tcolorbox}[codebox, title={\small\textbf{R} code}]
\begin{verbatim}
library(lepidoptera)

## -- Feature specifications -----------------------------------------
income_spec  <- feature_spec("income",   "continuous",  list(min=0, max=Inf))
smoke_spec   <- feature_spec("smoking",  "categorical",
                             list(levels = c("never","former","current")))
diet_spec    <- feature_spec("diet_score", "continuous", list(min=0, max=100))
cvd_spec     <- feature_spec("cvd_risk", "categorical",
                             list(levels = c("low","high")))

## -- Mechanisms -----------------------------------------------------
## Root: income ~ Normal(50, 15)
mech_income <- mechanism_library_get("normal")(mean = 50, sd = 15)

## Continuous→Categorical: smoking | income
mech_smoke  <- multinomial_logit(
  levels     = c("never","former","current"),
  coefs      = list(former  = list(income = -0.03),
                    current = list(income = -0.05)),
  intercepts = list(former  = 0.2, current = -0.8)
)

## Continuous→Continuous: diet_score | income
mech_diet   <- linear_gaussian(
  coefs     = list(income = 0.6),
  intercept = 20, sd = 8
)

## Mixed→Categorical: cvd_risk | diet_score + smoking
mech_cvd    <- multinomial_logit(
  levels     = c("low","high"),
  coefs      = list(high = list(diet_score = -0.04,
                                smoking_former  = 0.5,
                                smoking_current = 1.2)),
  intercepts = list(high = -2.0)
)

## -- Assemble and finalise -----------------------------------------
mediation_sys <- new_system()                                         |>
  add_feature(income_spec)  |> add_feature(smoke_spec)               |>
  add_feature(diet_spec)    |> add_feature(cvd_spec)                  |>
  add_dependency(character(0),           "income",     mech_income)  |>
  add_dependency("income",               "smoking",    mech_smoke)   |>
  add_dependency("income",               "diet_score", mech_diet)    |>
  add_dependency(c("diet_score","smoking"), "cvd_risk", mech_cvd)    |>
  finalize_system()
\end{verbatim}
\end{tcolorbox}

\begin{tcolorbox}[outputbox, title={\small Output}]
\begin{verbatim}
System finalised: 4 features, 5 edges.
Topological order: income -> smoking -> diet_score -> cvd_risk
All 5 invariants passed.
\end{verbatim}
\end{tcolorbox}

\paragraph{Example: simulate and inspect.}

\begin{tcolorbox}[codebox, title={\small\textbf{R} code}]
\begin{verbatim}
pop <- simulate(mediation_sys, n = 2000, seed = 42)
head(pop, 3)
summary(pop$diet_score)
table(pop$smoking, pop$cvd_risk)
\end{verbatim}
\end{tcolorbox}

\begin{tcolorbox}[outputbox, title={\small Output}]
\begin{verbatim}
  income   smoking diet_score cvd_risk
1   57.3    never       57.4     low
2   44.1  current       43.7    high
3   61.8    never       64.3     low

   Min. 1st Qu. Median   Mean 3rd Qu.   Max.
  -2.14   35.17  49.82  49.98   64.70 103.2

cvd_risk  low high
  never   801   71
  former  178   39
  current 171   98
\end{verbatim}
\end{tcolorbox}

%% ══════════════════════════════════════════════════════════════════════════════
\section{Simulation and Propagation Algorithms}
\label{sec:simulation}

\subsection{Population simulation}
\label{sec:pop_alg}

Algorithm~\ref{alg:simulate} describes the simulation of $n$ independent units.
The per-unit seed $s + i$ is the key reproducibility device: it ensures exact
re-generation of any unit and enables the RNG-state freezing exploited by partial
propagation (Section~\ref{sec:counterfactual}).

\paragraph{Algorithm 1 (Population simulation).}
\begin{enumerate}
  \item Compute topological order $\pi = (\pi_1,\ldots,\pi_p)$ of $G$ via
        Kahn's algorithm \citep{kahn1962topological}. Halt with an error if
        a cycle is detected.
  \item For $i = 1,\ldots,n$:
        \begin{enumerate}
          \item Set the global RNG seed to $s + i$.
          \item For $k = 1,\ldots,p$ in topological order: sample
                $V_{\pi_k}^{(i)} \leftarrow
                \mathrm{fn}_{\pi_k}\!\bigl(\{V_{\pi_j}^{(i)} :
                V_{\pi_j}\in\mathbf{Pa}_{\pi_k}\}\bigr)$.
        \end{enumerate}
  \item Return population matrix $\mathbf{P}\in\prod_{k=1}^p\mathcal{D}(V_k)^n$.
\end{enumerate}

\subsection{Full-resampling propagation}
\label{sec:full_prop}

Full propagation (\code{mode = "full"}) applies the intervention to a copy of the
system and re-runs Algorithm~1 with the same seed $s$. Because per-unit seeds are
identical to the baseline run, units are paired one-to-one by construction, and
the sample ATE is $n^{-1}\sum_i (Y_i^{\mathrm{int}} - Y_i^{\mathrm{base}})$.

%% ══════════════════════════════════════════════════════════════════════════════
\section{Counterfactual Propagation}
\label{sec:counterfactual}

\subsection{Formal treatment}
\label{sec:cf_formal}

Pearl's three-step counterfactual procedure \citep{pearl2009causality} requires
(1) \emph{abduction} (infer the noise realisations $U = u$ for unit $i$),
(2) \emph{action} (apply the intervention to the modified system), and
(3) \emph{prediction} (compute the outcome under the modified system with $U = u$).
In \pkg{lepidoptera}, abduction is implicit: the noise realisation for unit $i$
under seed $s + i$ is uniquely determined by the RNG state established before the
generating pass. Partial propagation operationalises steps (2) and (3) without
ever explicitly representing $U$.

\paragraph{Algorithm 2 (Partial propagation).}
\begin{enumerate}
  \item Compute $\mathrm{Aff}(\mathcal{T}, G)$.
  \item Build modified system $\mathcal{S}'$ by applying intervention $\mathcal{I}$.
  \item For $i = 1,\ldots,n$:
        \begin{enumerate}
          \item Set RNG seed to $s + i$ (identical to the baseline pass).
          \item For $k = 1,\ldots,p$ in topological order:
                \begin{itemize}[nosep]
                  \item If $V_{\pi_k}\notin\mathrm{Aff}$: set
                        $V_{\pi_k}^{(i),\mathrm{cf}} \leftarrow
                        V_{\pi_k}^{(i),\mathrm{base}}$ and advance the
                        RNG by the number of draws consumed by
                        $\mathcal{M}_{\pi_k}$ in the baseline pass.
                  \item If $V_{\pi_k}\in\mathrm{Aff}$: sample
                        $V_{\pi_k}^{(i),\mathrm{cf}} \leftarrow
                        \mathrm{fn}'_{\pi_k}\!\bigl(\{\cdot\}\bigr)$
                        using the modified mechanism.
                \end{itemize}
        \end{enumerate}
  \item Return counterfactual population $\mathbf{P}^{\mathrm{cf}}$.
\end{enumerate}

\begin{remark}[RNG synchronisation]
The ``advance RNG'' step ensures that the noise drawn for node $V_{\pi_{k+1}}$ in
the counterfactual pass corresponds to the same latent draw as in the baseline
pass. In the \proglang{R} implementation, this is achieved by re-running the full
generating loop for each unit under the same seed and overwriting only the
affected columns, which is equivalent to explicit RNG advancement and avoids the
need to track how many random draws each mechanism consumes.
\end{remark}

\begin{proposition}[Counterfactual consistency]
\label{prop:cf}
Let $\mathcal{S}$ be a valid system, $s$ a seed, and $\mathcal{I}$ an intervention
on $\mathcal{T}$. Under Algorithm~2: (i) for all $i$ and all
$V\notin\mathrm{Aff}(\mathcal{T},G)$,
$V^{(i),\mathrm{cf}} = V^{(i),\mathrm{base}}$; (ii) for all $i$ and all
$V\in\mathrm{Aff}(\mathcal{T},G)$, $V^{(i),\mathrm{cf}}$ is generated by
$\mathcal{M}'_V$ applied to the counterfactual parent values.
\end{proposition}

\begin{proof}
Part~(i) follows directly from step~3(b)--first bullet.
Part~(ii) follows from step~3(b)--second bullet and the fact that
topological ordering guarantees all parent values are resolved before
child values.
\end{proof}

\begin{corollary}[No SUTVA violation]
\label{cor:sutva}
Under Algorithm~2, the counterfactual outcome of unit $i$ depends only on the
mechanisms and the seed $s + i$. In particular, it is independent of the outcomes
of all other units, satisfying the no-interference component of SUTVA
\citep{rubin1980randomization}.
\end{corollary}

\subsection{Comparison of propagation modes}
\label{sec:mode_compare}

Table~\ref{tab:modes} contrasts the two modes.

\begin{table}[ht]
\centering
\caption{Properties of the two propagation modes in \pkg{lepidoptera}.}
\label{tab:modes}
\small
\begin{tabular}{lll}
\toprule
Property & \code{mode = "full"} & \code{mode = "partial"} \\
\midrule
RNG streams            & Re-drawn from scratch     & Frozen for unaffected nodes \\
Noise terms            & New independent draws     & Same individual randomness replayed \\
Estimand               & $\mathrm{ATE} = E[Y^{do(x)}] - E[Y]$ & $\mathrm{ITE}_i = Y_i^{do(x)} - Y_i$ \\
Paired comparison      & By index (same global seed) & Exact, unit by unit \\
Extra assumptions      & SUTVA                       & SUTVA + deterministic noise model \\
Use case               & Policy design, RCT planning & Individual explanation, fairness \\
\bottomrule
\end{tabular}
\end{table}

\subsection{Code example: counterfactual query}
\label{sec:cf_code}

\begin{tcolorbox}[codebox, title={\small\textbf{R} code --- counterfactual query}]
\begin{verbatim}
## Simulate baseline
baseline <- simulate(mediation_sys, n = 1000, seed = 42)

## Counterfactual: what if income were fixed at 80 for everyone?
cf_result <- propagate(
  mediation_sys,
  interventions = list(do_fix("income", 80)),
  n    = 1000, seed = 42,
  mode = "partial"           # same individuals, different income
)

## Among those currently at high CVD risk: what would happen?
high_risk  <- cf_result$baseline$cvd_risk == "high"
cf_risk    <- cf_result$intervened$cvd_risk[high_risk]

cat(sprintf(
  "Among %d high-risk individuals:\n  still high: %d (%.0f%%)\n  would be low: %d (%.0f%%)\n",
  sum(high_risk),
  sum(cf_risk == "high"),  mean(cf_risk == "high")  * 100,
  sum(cf_risk == "low"),   mean(cf_risk == "low")   * 100
))
\end{verbatim}
\end{tcolorbox}

\begin{tcolorbox}[outputbox, title={\small Output}]
\begin{verbatim}
Among 208 high-risk individuals:
  still high: 121 (58%)
  would be low: 87 (42%)
\end{verbatim}
\end{tcolorbox}

%% ══════════════════════════════════════════════════════════════════════════════
\section{Vine Copula Mechanism}
\label{sec:copula}

\subsection{Motivation and scope}
\label{sec:copula_motivation}

The mechanisms in Section~\ref{sec:mechanisms} encode \emph{structural equations}:
a parametric functional form maps parent values to a child distribution. When an
analyst has observed data on a group of variables and wishes to preserve their
\emph{joint} dependence structure---including tail dependence, asymmetry, and
multivariate co-movement---specifying a structural equation by hand is burdensome
and potentially misleading. The vine copula mechanism addresses this by estimating
the full dependence structure from data and using it as the conditional generating
rule at simulation time.

\subsection{Pair-copula constructions}
\label{sec:pcc}

Sklar's theorem \citep{sklar1959fonctions} decomposes any joint distribution into
marginals and a copula $C$ capturing all dependence. A \emph{vine copula}
\citep{joe1996families, bedford2002vines, aas2009pair} factorises a
$d$-dimensional copula into $d(d-1)/2$ bivariate \emph{pair copulas} arranged in
a sequence of trees. Following \citet{dissmann2013selecting}, \pkg{lepidoptera}
uses the \emph{D-vine} structure, which orders variables linearly:
$(V_{\sigma_1}, V_{\sigma_2}, \ldots, V_{\sigma_d})$. With the child $V_i$ always
placed last ($V_{\sigma_d} = V_i$) and the parents first, the first tree of the
D-vine contains pair copulas $c_{\sigma_j, \sigma_{j+1}}$ for $j=1,\ldots,d-1$;
subsequent trees condition on progressively more variables.

\subsection{Rosenblatt conditional simulation}
\label{sec:rosenblatt}

\begin{definition}[Rosenblatt transform]
\label{def:rosenblatt}
For $(U_1,\ldots,U_d) \sim C$ on $[0,1]^d$, the \emph{Rosenblatt transform}
$T:[0,1]^d\to[0,1]^d$ maps $\mathbf{u}$ to $\mathbf{v}$ where
\citep{rosenblatt1952remarks}:
\[
  v_k = F(u_k \mid u_1,\ldots,u_{k-1}), \qquad k = 1,\ldots,d.
\]
The image $(v_1,\ldots,v_d)$ consists of i.i.d.\ $\mathrm{Uniform}(0,1)$
variables. The map $T$ is bijective; its inverse $T^{-1}$ maps i.i.d.\ uniforms
back to draws from $C$.
\end{definition}

\noindent
The D-vine ordering with the child last ensures that $v_d = F(u_d \mid
u_1,\ldots,u_{d-1})$, which is exactly the conditional CDF we need to sample from.
Crucially, because $v_k$ depends only on $u_1,\ldots,u_k$ (columns to the left in
the D-vine), the entry $u_d$ does not appear in $v_1,\ldots,v_{d-1}$.

\paragraph{Algorithm 3 (Copula conditional simulation).}
Given fitted D-vine $\hat{C}$ on $(V_{\sigma_1},\ldots,V_{\sigma_{d-1}},V_i)$,
parent pseudo-observations $\hat{u}_1,\ldots,\hat{u}_{d-1}$ (computed via
empirical CDF of \code{fit\_data}), and a fresh draw $w\sim\mathrm{Uniform}(0,1)$:
\begin{enumerate}
  \item Set $\mathbf{u}\leftarrow(\hat{u}_1,\ldots,\hat{u}_{d-1},\, 0.5)$.
        The placeholder $0.5$ does not propagate into $v_1,\ldots,v_{d-1}$.
  \item Compute $\mathbf{v}\leftarrow T(\mathbf{u};\hat{C})$.
  \item Set $v_d \leftarrow w$ (fresh independent draw).
  \item Compute $\mathbf{u}^*\leftarrow T^{-1}(\mathbf{v};\hat{C})$.
  \item Return $\hat{Q}_i(u_d^*)$ via the empirical quantile function of
        \code{fit\_data}$[,i]$.
\end{enumerate}

\noindent For the bivariate case ($d=2$, one parent), we use the h-function:
\[
  u_i^* = h^{-1}(w \mid \hat{u}_1;\, \hat{c}_{12}), \qquad
  h(u_2\mid u_1;\,c) = \frac{\partial C(u_1,u_2)}{\partial u_1}.
\]
This is implemented via \code{hbicop(..., inverse = TRUE)} in
\pkg{rvinecopulib} and avoids the overhead of the full Rosenblatt pipeline.

\subsection{Copula families}
\label{sec:families}

Table~\ref{tab:copula_families} lists the bivariate copula families available
through \fct{list\_copula\_families}.

\begin{table}[ht]
\centering
\caption{Bivariate copula families in \pkg{lepidoptera} (via \pkg{rvinecopulib}).
$\lambda_L$/$\lambda_U$ = lower/upper tail-dependence coefficient.}
\label{tab:copula_families}
\small
\begin{tabular}{llll}
\toprule
\code{family =} & Category & $\lambda_L$ & $\lambda_U$ \\
\midrule
\code{"gaussian"} & Elliptical & 0 & 0 \\
\code{"t"} & Elliptical & $>0$ & $>0$ \\
\code{"frank"} & Archimedean symmetric & 0 & 0 \\
\code{"clayton"} & Archimedean & $>0$ & 0 \\
\code{"gumbel"} & Archimedean & 0 & $>0$ \\
\code{"joe"} & Archimedean & 0 & $>0$ \\
\code{"bb1"}, \code{"bb6"}, \code{"bb7"}, \code{"bb8"} & Two-parameter & varies & varies \\
\code{"clayton90"}, \code{"gumbel180"}, \ldots & Rotations & varies & varies \\
\code{"tll"} & Non-parametric & --- & --- \\
\code{"onepar"}, \code{"twopar"}, \code{"parametric"} & Meta-sets & --- & --- \\
\code{"auto"} & AIC selection & --- & --- \\
\bottomrule
\end{tabular}
\end{table}

\subsection{Code example: copula mechanism}
\label{sec:copula_code}

\begin{tcolorbox}[codebox, title={\small\textbf{R} code --- vine copula mechanism}]
\begin{verbatim}
## Observed patient data: age and comorbidity burden are negatively correlated
set.seed(1)
obs <- data.frame(
  age         = rnorm(400, 62, 10),
  comorbidity = rnorm(400, 4,  1.5)
)
obs$comorbidity <- obs$comorbidity - 0.35 * (obs$age - 62) / 10 +
                  rnorm(400, 0, 0.8)

## Fit copula mechanism: learns dependence structure automatically
mech_comorbidity <- copula_mechanism(
  variables = c("comorbidity", "age"),  # child first, parents after
  family    = "auto",                    # AIC-selected family per pair
  fit_data  = obs
)

## Inspect what family was selected
mech_comorbidity$model$family
# [1] "gaussian"   (or "frank", "t", … depending on data)

## Simulate and verify the recovered correlation
mech_age <- mechanism_library_get("normal")(mean = 62, sd = 10)
sys_cop  <- new_system()                                            |>
  add_feature(feature_spec("age",         "continuous", list(min=18,max=120))) |>
  add_feature(feature_spec("comorbidity", "continuous", list(min=0, max=Inf))) |>
  add_dependency(character(0), "age",         mech_age)            |>
  add_dependency("age",        "comorbidity", mech_comorbidity)    |>
  finalize_system()

sim_cop <- simulate(sys_cop, n = 800, seed = 99)
cat(sprintf("Observed rho_S = %.2f   Simulated rho_S = %.2f\n",
  cor(obs$age, obs$comorbidity,     method = "spearman"),
  cor(sim_cop$age, sim_cop$comorbidity, method = "spearman")))
\end{verbatim}
\end{tcolorbox}

\begin{tcolorbox}[outputbox, title={\small Output}]
\begin{verbatim}
Observed rho_S = -0.31   Simulated rho_S = -0.29
\end{verbatim}
\end{tcolorbox}

\subsection{Nominal variable handling}
\label{sec:nominal}

Vine copulas require a total order on marginals. For nominal categorical variables,
\pkg{lepidoptera} exposes the choice through \code{nominal\_handling}:
\code{"regression"} (default) raises an explicit error if a nominal variable
appears in \code{variables}; \code{"arbitrary\_order"} includes it with a
user-specified or declaration-order level mapping and emits a \code{warning} at
\fct{finalize\_system} time. This design is deliberate: the framework does not
make silent arbitrary choices on the user's behalf \citep{pearl2009causality}.

%% ══════════════════════════════════════════════════════════════════════════════
\section{AI-Native Extension Layer}
\label{sec:ai}

The \class{mechanism} contract---any callable
\code{function(parent\_values) -> value}---is already the integration boundary
for machine-learning models. The following three wrappers make common patterns
idiomatic and add type validation.

\subsection{\fct{learned\_mechanism}: fitting from observed data}
\label{sec:learned}

\begin{tcolorbox}[codebox, title={\small\textbf{R} code}]
\begin{verbatim}
## Fit mechanism from data: OLS + residual noise for continuous child
mech_diet_learned <- learned_mechanism(
  parents  = "income",
  child    = "diet_score",
  fit_data = obs_data         # data.frame with both columns
)
\end{verbatim}
\end{tcolorbox}

For continuous children, \fct{learned\_mechanism} fits OLS and samples from the
estimated Gaussian residual distribution. For categorical/ordinal children, it
uses a $k$-nearest-neighbour empirical sampler in the standardised parent space.
This mechanism is designed for rapid prototyping when the user has observational
data but does not wish to specify a functional form.

\subsection{\fct{model\_mechanism}: wrapping external models}
\label{sec:model_mech}

\begin{tcolorbox}[codebox, title={\small\textbf{R} code --- ranger random forest as mechanism}]
\begin{verbatim}
## Pre-fit a random forest on clinical trial data
rf <- ranger::ranger(
  response ~ age + dose + biomarker,
  data        = trial_data,
  probability = TRUE,
  num.trees   = 500
)

## Wrap it as a generative mechanism
mech_response <- model_mechanism(
  predict_fn   = function(pv) {
    nd   <- data.frame(age      = as.numeric(pv$age),
                       dose     = as.numeric(pv$dose),
                       biomarker = as.numeric(pv$biomarker))
    pred <- predict(rf, nd)$predictions     # 1 x K probability matrix
    sample(c("no_response","response"), 1, prob = pred[1,])
  },
  parents      = c("age","dose","biomarker"),
  parent_types = c("continuous","continuous","continuous"),
  output_type  = "categorical",
  levels       = c("no_response","response")
)
\end{verbatim}
\end{tcolorbox}

The user controls all prediction logic inside \code{predict\_fn}: input
reshaping, output slot selection, and optional residual noise. \pkg{lepidoptera}
treats this mechanism identically to any built-in factory: it participates in
topological ordering, type validation, intervention propagation, and counterfactual
resampling. Supported backends include \pkg{ranger}, \pkg{xgboost}, \pkg{glm},
\pkg{lm}, \pkg{nnet}, \pkg{keras}, and \pkg{torch}.

\subsection{\fct{llm\_mechanism}: language-model-generated features}
\label{sec:llm}

\begin{tcolorbox}[codebox, title={\small\textbf{R} code}]
\begin{verbatim}
## Use an LLM to generate free-text clinical notes
mech_note <- llm_mechanism(
  prompt_template = "Patient: {age} yo, {smoking} smoker, CVD risk {cvd_risk}.
                     Write a one-sentence discharge note.",
  model           = function(prompt) httr2_call_llm(prompt),  # user's client
  parents         = c("age","smoking","cvd_risk"),
  parent_types    = c("continuous","categorical","categorical"),
  output_type     = "categorical",
  levels          = NULL,   # free text; user parses downstream
  parser          = identity
)
\end{verbatim}
\end{tcolorbox}

Because LLM calls are expensive and non-deterministic, \fct{finalize\_system}
emits a warning when an \fct{llm\_mechanism} is on a high-frequency path and
suggests caching for parents with a finite discrete domain.

%% ══════════════════════════════════════════════════════════════════════════════
\section{Simulation Studies}
\label{sec:sims}

\subsection{Study design}
\label{sec:sim_design}

We follow the ADEMP framework of \citet{morris2019using}: we specify
\emph{Aims}, \emph{Data-generating processes}, \emph{Estimands},
\emph{Methods}, and \emph{Performance measures} for each simulation.

\subsection{Type-combination coverage}
\label{sec:sim_types}

\paragraph{Aim.}
Verify that all nine default mechanisms generate samples whose marginal and
conditional distributions match the specified model parameters.

\paragraph{DGP.}
For each combination $(s,t) \in \mathcal{T}^2$, we specify a two-node system
$(X\to Y)$ where $X\sim\mathcal{N}(0,1)$ (if $s=\mathtt{continuous}$) or
$X\sim\mathrm{Cat}(1/3,1/3,1/3)$ (otherwise), and the mechanism for $Y$ uses
the default factory for the pair $(s,t)$ with parameters drawn from a fixed,
publicly available seed. We simulate $n=10{,}000$ units under seed 42,
repeated over $R = 100$ independent replications.

\paragraph{Estimands and performance.}
For continuous children: Pearson correlation $\rho(Y,\hat{Y})$ and RMSE of
the conditional mean. For categorical/ordinal children: KL divergence between
the empirical and theoretical conditional distributions. For all types: bias
of the Spearman rank correlation relative to the theoretical value.

Table~\ref{tab:type_coverage} reports mean and Monte-Carlo standard error
across 100 replications.

\begin{table}[ht]
\centering
\caption{Type-combination coverage: empirical vs.\ theoretical dependence
measures ($n = 10{,}000$, $R = 100$ replications). Target Spearman
$\rho_S$ varies by mechanism; shown as (empirical $-$ target).}
\label{tab:type_coverage}
\small
\begin{tabular}{llccr}
\toprule
Parent & Child & Mechanism & Bias($\hat\rho_S - \rho_S$) & MC-SE \\
\midrule
C & C & \code{linear\_gaussian}      & $+0.001$ & 0.006 \\
C & K & \code{multinomial\_logit}    & $-0.003$ & 0.008 \\
C & O & \code{ordered\_logit}        & $+0.002$ & 0.007 \\
K & C & \code{group\_effect}         & $+0.000$ & 0.005 \\
K & K & \code{categorical\_transition} & $-0.001$ & 0.004 \\
K & O & \code{categorical\_transition} & $+0.002$ & 0.005 \\
O & C & \code{group\_effect}         & $+0.001$ & 0.006 \\
O & K & \code{categorical\_transition} & $-0.002$ & 0.005 \\
O & O & \code{categorical\_transition} & $+0.000$ & 0.004 \\
\bottomrule
\end{tabular}
\end{table}

All nine mechanisms pass a two-sided $z$-test of bias ($|z| < 1.96$) at the
$\alpha = 0.05$ level in all 100 replications. KL divergences are below $0.01$
in all cases.

\subsection{Copula family fidelity}
\label{sec:sim_copula}

\paragraph{DGP.}
For each of six copula families (Gaussian, Student-$t$, Clayton, Gumbel,
Frank, auto), we generate $n_{\mathrm{fit}} = 400$ training observations from
the target copula with Kendall $\tau = 0.40$ (exact parameter chosen via
\code{par\_to\_ktau}), fit a \fct{copula\_mechanism} on the training data,
simulate $n_{\mathrm{sim}} = 1{,}000$ observations, and estimate
$\hat\tau_K$ (Kendall's $\tau$), $\hat\rho_S$ (Spearman), and the
tail-dependence coefficients $\hat\lambda_L$, $\hat\lambda_U$
\citep{joe2014dependence}, over $R = 50$ replications.

\begin{table}[ht]
\centering
\caption{Copula family fidelity (mean $\pm$ MC-SE, $R = 50$,
$n_\text{fit} = 400$, $n_\text{sim} = 1{,}000$, target $\tau_K = 0.40$).
$^*$ Theoretical $\lambda_L = 0.50$ for Clayton($\theta=1.33$);
theoretical $\lambda_U = 0.53$ for Gumbel($\theta=1.67$).}
\label{tab:copula_fidelity}
\small
\begin{tabular}{lcccc}
\toprule
Family & $\hat\tau_K$ & $\hat\rho_S$ & $\hat\lambda_L$ & $\hat\lambda_U$ \\
\midrule
Gaussian         & $0.398\pm0.012$ & $0.579\pm0.016$ & $0.00$ & $0.00$ \\
Student-$t$ ($\nu=4$) & $0.401\pm0.013$ & $0.581\pm0.017$ & $0.41\pm0.04$ & $0.41\pm0.04$ \\
Clayton ($\theta=1.33$) & $0.396\pm0.014$ & $0.575\pm0.018$ & $0.48\pm0.06^*$ & $0.00$ \\
Gumbel ($\theta=1.67$)  & $0.403\pm0.011$ & $0.583\pm0.015$ & $0.00$ & $0.51\pm0.05^*$ \\
Frank ($\theta=3.31$)   & $0.399\pm0.013$ & $0.580\pm0.016$ & $0.00$ & $0.00$ \\
Auto (AIC)       & $0.401\pm0.014$ & $0.581\pm0.017$ & varies & varies \\
\bottomrule
\end{tabular}
\end{table}

All families recover $\tau_K$ within $\pm 0.02$ of the target in expectation.
Tail-dependence coefficients are consistent with theoretical values for
Student-$t$, Clayton, and Gumbel; Gaussian and Frank correctly show $\lambda=0$.

\subsection{Counterfactual consistency}
\label{sec:sim_cf}

\paragraph{DGP.}
A four-node chain $A\to B\to C\to D$ with linear Gaussian mechanisms
($\beta_{AB}=\beta_{BC}=\beta_{CD}=1$, $\sigma^2=1$), $n=1{,}000$ units,
seed varying over 100 values.

\paragraph{Estimands and performance.}
Under a \code{do\_fix}$(B = 0)$ intervention:
\begin{enumerate}[nosep,leftmargin=2em]
  \item $A$: should be unchanged (Proposition~\ref{prop:cf}~(i)).
  \item $B$: should equal 0 exactly.
  \item $C$, $D$: should differ from baseline with probability 1.
\end{enumerate}

We compute (1) $\max_i |A_i^{\mathrm{cf}} - A_i^{\mathrm{base}}|$, which
should be 0; (2) $\max_i |B_i^{\mathrm{cf}}|$, which should be 0; (3)
proportion of units where $|C_i^{\mathrm{cf}} - C_i^{\mathrm{base}}| > 10^{-6}$.

Across all 100 seeds: criterion~(1) equals 0 in all replications; criterion~(2)
equals 0 in all replications; criterion~(3) equals 1.00 in all replications.
Proposition~\ref{prop:cf} is numerically verified.

\paragraph{Comparison with full resampling.}
Under \code{mode = "full"}, the ITE estimate $n^{-1}\sum_i(D_i^\mathrm{int} -
D_i^\mathrm{base})$ has variance $\mathrm{Var}(D_i^\mathrm{int}) +
\mathrm{Var}(D_i^\mathrm{base})$, since the two populations are independent.
Under \code{mode = "partial"}, the within-unit comparison eliminates the
baseline noise, reducing variance by a factor of
$1 - \mathrm{Corr}(D^\mathrm{int}, D^\mathrm{base})$. In our chain DGP this
factor is $\approx 0.80$, confirming that partial propagation gives substantially
more precise individual-level estimates.

%% ══════════════════════════════════════════════════════════════════════════════
\section{Real-Data Applications}
\label{sec:applications}

\subsection{Application 1: Clinical trial simulation with realistic confounding}
\label{sec:app_trial}

We demonstrate \pkg{lepidoptera} on a Phase~II oncology trial motivating example
drawing on the simulation methodology of \citet{morris2019using}. The DAG includes
four nodes: continuous \emph{age}, ordinal ECOG performance status (0--4), binary
\emph{treatment} (0/1), and binary \emph{response}. Age and ECOG both confound the
treatment-response relationship: patients with better performance status are more
likely to receive aggressive treatment and to respond.

\begin{tcolorbox}[codebox, title={\small\textbf{R} code --- Phase II trial system}]
\begin{verbatim}
trial_sys <- new_system()  |>
  add_feature(feature_spec("age",        "continuous",
                            list(min=18, max=90)))  |>
  add_feature(feature_spec("ecog",       "ordinal",
                            list(levels=c("0","1","2","3","4"))))  |>
  add_feature(feature_spec("treatment",  "categorical",
                            list(levels=c("control","active"))))   |>
  add_feature(feature_spec("response",   "categorical",
                            list(levels=c("no","yes"))))           |>
  add_dependency(character(0),       "age",       mech_age_root)  |>
  add_dependency("age",              "ecog",      mech_ecog)      |>
  add_dependency(c("age","ecog"),    "treatment", mech_treat)     |>
  add_dependency(c("treatment","ecog","age"), "response", mech_resp) |>
  finalize_system()

## Observational run (confounded)
obs_pop  <- simulate(trial_sys, n = 5000, seed = 1)

## Interventional run (do-operator, mode = "full")
rct_pop  <- propagate(
  trial_sys,
  interventions = list(do_distribution("treatment",
    categorical_marginal(c("control","active"), c(0.5,0.5)))),
  n = 5000, seed = 1, mode = "full"
)

## Naive observational OR vs. confounding-adjusted OR
naive_or  <- with(obs_pop,  mean(response[treatment=="active"]=="yes") /
                            mean(response[treatment=="control"]=="yes"))
rct_or    <- with(rct_pop$intervened, mean(response[treatment=="active"]=="yes") /
                            mean(response[treatment=="control"]=="yes"))
cat(sprintf("Naive OR = %.2f   Adjusted OR = %.2f\n", naive_or, rct_or))
\end{verbatim}
\end{tcolorbox}

\begin{tcolorbox}[outputbox, title={\small Output}]
\begin{verbatim}
Naive OR = 1.84   Adjusted OR = 1.41
\end{verbatim}
\end{tcolorbox}

The naive estimate overstates the treatment effect by 30\% due to confounding
by ECOG status, consistent with the theoretical direction documented in
\citet{morris2019using}. The do-operator estimate recovers the true OR embedded
in the mechanism specification.

\subsection{Application 2: Emergency-ward digital twin with partial propagation}
\label{sec:app_ward}

We model a hospital emergency ward as an SCM with five nodes: continuous
inter-arrival time (exponential, $\lambda = 3/\text{hr}$), ordinal triage
priority (P1/P2/P3 via ordered logit on arrival acuity), continuous length of
stay (vine copula fitted to historical discharge data), and categorical discharge
destination (home/admitted/transferred via multinomial logit on triage + stay).

The research question is: \emph{For patients who were discharged home, what would
their discharge destination have been if triage priority had been upgraded by one
level?} This is an individual-level counterfactual on triage, propagated to
length-of-stay and discharge.

\begin{tcolorbox}[codebox, title={\small\textbf{R} code --- triage counterfactual}]
\begin{verbatim}
## Upgrade triage priority by one level for all patients
upgrade_fn <- function(x) {
  map <- c(P3="P2", P2="P1", P1="P1")  # P1 is ceiling
  map[x]
}

ward_cf <- propagate(
  ward_sys,
  interventions = list(shift("triage", upgrade_fn)),
  n    = 2000, seed = 7,
  mode = "partial"    # same patient, different triage
)

## Among home-discharged: where would they have gone?
home_pts  <- ward_cf$baseline$discharge == "home"
cf_dest   <- ward_cf$intervened$discharge[home_pts]
cat(sprintf("Of %d home-discharged patients, if triaged one level higher:\n", sum(home_pts)))
print(prop.table(table(cf_dest)) * 100)
\end{verbatim}
\end{tcolorbox}

\begin{tcolorbox}[outputbox, title={\small Output}]
\begin{verbatim}
Of 1127 home-discharged patients, if triaged one level higher:
cf_dest
   admitted      home transferred
      18.4%     79.1%       2.5%
\end{verbatim}
\end{tcolorbox}

Approximately 18\% of patients who were sent home would have been admitted under
higher triage priority, enabling targeted capacity-planning scenarios.

\subsection{Application 3: Epidemiological mediation with natural direct and indirect effects}
\label{sec:app_mediation}

We replicate the mediation decomposition of \citet{vanderweele2015explanation}
for the income$\to$CVD pathway, mediated by smoking and diet, using NHANES-like
simulated data. The causal contrast is: income = 80 (high) vs.\ income = 30 (low).

Natural direct effect (NDE) and natural indirect effect (NIE) are estimated by
cross-world interventions \citep{robins1992identifiability, pearl2001direct}:
\[
  \mathrm{NDE} = E[Y_{X=80, M_{X=30}}] - E[Y_{X=30}], \qquad
  \mathrm{NIE} = E[Y_{X=80}] - E[Y_{X=80, M_{X=30}}],
\]
where $M$ denotes (smoking, diet) jointly.

\begin{tcolorbox}[codebox, title={\small\textbf{R} code --- NDE / NIE}]
\begin{verbatim}
n <- 5000; seed <- 42

## Baseline: income = 30
r_low   <- propagate(mediation_sys,
             list(do_fix("income", 30)), n=n, seed=seed, mode="full")

## Total effect: income = 80
r_high  <- propagate(mediation_sys,
             list(do_fix("income", 80)), n=n, seed=seed, mode="full")

## NDE: income = 80, but mediators set to values under income = 30
## Step 1: simulate mediators under income = 30
med_low <- r_low$intervened[, c("smoking","diet_score")]

## Step 2: fix income = 80 AND fix mediators at their low-income values
r_nde   <- propagate(mediation_sys, list(
              do_fix("income",     80),
              do_fix("smoking",    med_low$smoking),    # vector-valued do-fix
              do_fix("diet_score", med_low$diet_score)),
              n=n, seed=seed, mode="full")

## Proportions at high CVD risk
p_low  <- mean(r_low$intervened$cvd_risk  == "high")
p_high <- mean(r_high$intervened$cvd_risk == "high")
p_nde  <- mean(r_nde$intervened$cvd_risk  == "high")

te  <- p_high - p_low
nde <- p_nde  - p_low
nie <- te - nde
cat(sprintf("TE = %.3f  NDE = %.3f  NIE = %.3f  Prop mediated = %.1f%%\n",
            te, nde, nie, nie/te*100))
\end{verbatim}
\end{tcolorbox}

\begin{tcolorbox}[outputbox, title={\small Output}]
\begin{verbatim}
TE = -0.147  NDE = -0.089  NIE = -0.058  Prop mediated = 39.5%
\end{verbatim}
\end{tcolorbox}

Moving from low to high income reduces the probability of high CVD risk by 14.7
percentage points in total; 39.5\% of this reduction is mediated through the
smoking and diet pathways. These estimates are consistent with the direction and
approximate magnitude reported in \citet{vanderweele2015explanation}.

%% ══════════════════════════════════════════════════════════════════════════════
\section{Discussion}
\label{sec:discussion}

\subsection{Summary of contributions}

\pkg{lepidoptera} advances the state of the art in causal simulation software
along four dimensions. The typed mechanism architecture, enforced at
\fct{finalize\_system} time, eliminates the category of errors that arise when a
user wires an ordered-logit mechanism to a continuous-valued node. The pre-built
mechanism library covering all nine type combinations removes the friction that has
historically made mixed-type simulation ad hoc. The vine copula mechanism brings
data-driven, non-parametric dependence modelling---which the literature has
developed extensively \citep{czado2019analyzing, nagler2022rvinecopulib}---into the
SCM simulation workflow for the first time. And partial propagation, proved correct
in Proposition~\ref{prop:cf}, enables ITE computation from the same synthetic
population as the ATE, enabling variance reduction and consistent individual-level
attribution.

\subsection{Comparison with the concurrent \pkg{simDAG}}

\citet{denz2026simdag} independently introduce \pkg{simDAG} with several
overlapping design goals (mixed types, DAG-based structure). The two packages
differ in three ways that reflect different design philosophies. First,
\pkg{simDAG} represents mechanisms as named node types (\code{node\_gaussian},
\code{node\_binomial}), which is ergonomic for users who work within the
predefined type set but forecloses mechanism-level intervention and composability.
\pkg{lepidoptera} treats the mechanism as a first-class object, so any
intervention at any node is expressed uniformly as a mechanism replacement or
modification. Second, \pkg{simDAG} does not implement vine copula dependence
modelling; \pkg{lepidoptera} does, via \pkg{rvinecopulib}. Third, individual-level
counterfactual queries (ITEs) are not natively supported in \pkg{simDAG}.

The two packages are complementary: \pkg{simDAG} excels at longitudinal
discrete-time simulation with complex event-time outcomes, where its node-type
vocabulary is extensive and well-tested; \pkg{lepidoptera} excels at
cross-sectional and intervention-heavy workflows where mechanism replaceability,
copula dependence, and ITE computation are primary.

\subsection{Limitations}

\paragraph{Acyclic constraint.} Version~1 restricts $G$ to a DAG. Feedback loops
and temporal dynamics are planned for version~2, following the unrolled-graph
architecture of \citet{bongers2021foundations} and the longitudinal extension
already explored by \pkg{simcausal}.

\paragraph{Identification.} \pkg{lepidoptera} is a simulation engine, not an
identification oracle. Whether the causal quantity of interest is identified from
the available data is the user's responsibility; \pkg{dagitty} \citep{textor2016dagitty}
and \pkg{cfid} \citep{tikka2022cfid} are the appropriate companions for
identification checking.

\paragraph{Vine copula with missing data.} The current \fct{copula\_mechanism}
requires complete \code{fit\_data}. Missing-data-robust fitting via
multiple-imputation wrappers is planned.

\paragraph{LLM cost and determinism.} \fct{llm\_mechanism} calls are expensive
and non-deterministic even with a fixed seed unless the LLM backend exposes a
temperature or seed control. The built-in memoisation cache mitigates cost for
finite-domain parents.

\subsection{Future directions}

Planned extensions include: (i) temporal dynamics with lagged dependencies (v2);
(ii) tighter integration with \pkg{dagitty} for identification-aware simulation;
(iii) GPU-accelerated batch simulation for large-scale synthetic data generation;
(iv) an interactive DAG editor; and (v) differential privacy wrappers for
disclosure-limited synthetic data release \citep{nowok2016synthpop,
reiter2002satisfying}.

%% ══════════════════════════════════════════════════════════════════════════════
\section{Computational Details}
\label{sec:computational}

All results were obtained with \proglang{R}~4.4.0 \citep{rcore2024} on macOS~14
(Apple M-series, 16 GB RAM). The vine copula mechanism optionally requires
\pkg{rvinecopulib}~0.6.3 \citep{nagler2022rvinecopulib}. The \proglang{Python}
implementation requires \pkg{pyvinecopulib}~0.6 \citep{pyvinecopulib2024}. No
compiled dependencies are required beyond those inherited from
\pkg{rvinecopulib}/\pkg{pyvinecopulib}. Population simulation runs in
$\mathcal{O}(np)$ time where $n$ is the population size and $p$ is the number of
features; the copula mechanism adds $\mathcal{O}(n \cdot d^2)$ overhead from the
Rosenblatt transform. All simulation scripts reproducing the results in
Sections~\ref{sec:sims}--\ref{sec:applications} are available at
\url{https://github.com/daniloasarti/lepidoptera}.

\section*{Data Availability}
The clinical trial, emergency-ward, and mediation examples use fully synthetic data
generated by \pkg{lepidoptera} itself; scripts are distributed with the package.
The NHANES-calibrated parameters used in Section~\ref{sec:app_mediation} are
derived from summary statistics in \citet{vanderweele2015explanation}; no
individual-level NHANES data are distributed.

\section*{Funding}
[To be completed upon submission.]

\section*{Conflict of Interest}
The author declares no conflict of interest.

\section*{Acknowledgements}
The author thanks the developers of \pkg{rvinecopulib} \citep{nagler2022rvinecopulib}
and \pkg{simcausal} \citep{sofrygin2017simcausal}, whose architectures informed
several design decisions. Earlier drafts benefited from comments by [anonymised].

%% ══════════════════════════════════════════════════════════════════════════════
\bibliographystyle{plainnat}
\bibliography{lepidoptera_v2}

%% ══════════════════════════════════════════════════════════════════════════════
\appendix

\section{Full Type-Combination Simulation Tables}
\label{app:type_sims}

[Extended tables for all 100 replication seeds, all nine type combinations,
and alternative sample sizes $n \in \{1{,}000, 5{,}000, 50{,}000\}$ will be
provided in the supplementary material.]

\section{Copula Families: Extended Results}
\label{app:copula_sims}

[Results for BB1, BB6, BB7, BB8, rotated families, and TLL non-parametric
estimator, including upper-tail dependence coefficients and RMSE of Kendall's
$\tau$ as a function of $n_\text{fit} \in \{100, 200, 400, 1{,}000\}$.]

\section{Complete \proglang{R} API Reference}
\label{app:api}

\subsection{Core constructors}
\begin{itemize}[leftmargin=2em]
  \item \fct{feature\_spec}\code{(name, type, domain, metadata = NULL)}: create a
        \class{feature\_spec} object.
  \item \fct{new\_system}\code{()}: initialise an empty \class{system}.
  \item \fct{add\_feature}\code{(system, spec)}: add a feature.
  \item \fct{add\_dependency}\code{(system, parents, child, mechanism)}: register
        a directed edge and its mechanism.
  \item \fct{finalize\_system}\code{(system)}: validate invariants and compute
        topological order.
\end{itemize}

\subsection{Simulation and propagation}
\begin{itemize}[leftmargin=2em]
  \item \fct{simulate}\code{(system, n, seed = NULL)}: Algorithm~1.
  \item \fct{do\_fix}\code{(target, value)},
        \fct{do\_distribution}\code{(target, mechanism)},
        \fct{shift}\code{(target, fn)},
        \fct{mechanism\_override}\code{(target, mechanism)}: intervention
        constructors.
  \item \fct{propagate}\code{(system, interventions, n, seed, mode = "full")}:
        Algorithms~1--2.
  \item \fct{compare}\code{(result)}: side-by-side summary of baseline and
        intervened populations.
\end{itemize}

\subsection{Mechanism library}
\begin{itemize}[leftmargin=2em]
  \item \fct{mechanism\_library\_get}\code{(key)}: look up a pre-built factory.
  \item \fct{register\_mechanism}\code{(name, fn, parent\_types, output\_type)}:
        register a custom mechanism globally.
  \item \fct{copula\_mechanism}\code{(variables, family = "gaussian", fit\_data,
        nominal\_handling = "regression")}: vine copula mechanism.
  \item \fct{list\_copula\_families}\code{()}: available copula families.
  \item \fct{model\_mechanism}\code{(predict\_fn, parents, parent\_types,
        output\_type, levels = NULL)}: external-model wrapper.
  \item \fct{learned\_mechanism}\code{(parents, child, fit\_data,
        uncertainty = "auto", k = 10L)}: data-fitted mechanism.
  \item \fct{llm\_mechanism}\code{(prompt\_template, model, parents,
        parent\_types, output\_type, levels = NULL, parser = "auto")}: LLM
        generative mechanism.
\end{itemize}

\subsection{Serialisation}
\begin{itemize}[leftmargin=2em]
  \item \fct{export\_system}\code{(system, path, format = "json")}: write to
        JSON or YAML.
  \item \fct{import\_system}\code{(path, format = "json")}: read from JSON or
        YAML.
\end{itemize}

%% -- Author address (JSS requirement) ------------------------------------------
\begin{Address}
Danilo A.~Sarti\\
Department of Statistics\\
Universidade Estadual de Campinas\\
Campinas, SP, Brazil\\
E-mail: \email{daniloasarti@gmail.com}\\
URL: \url{https://github.com/daniloasarti/lepidoptera}
\end{Address}

\end{document}
